% 原文 pp.129–136（印刷 pp.1247–1254）
\section{Lecture 9：孤子}

\begin{keytakeaways}
本讲把「孤子」定义为庞加莱不变场论之经典场方程的行波解：空间局域、能量有限，并在给定同伦类中取能量的全局极小。核心数学工具是把孤子与「有限能量经典解空间的连通分支」联系起来：孤子在解空间中的 $P$-轨道 $\mathcal{O}=P/P_s$ 同构于余切丛 $T^*\mathcal{O}^+_m$，其量子化给出 $P$ 的不可约表示 $\mathbf{L}^2(\mathcal{O}^+_m)$，且孤子分支的 Hilbert 空间没有真空、能谱有正的下界。孤子质量随作用量整体常数 $C$ 线性增长（$\sim C$），远重于微扰粒子，故微扰论看不见孤子。在 $1+1$ 维含费米子的模型中，孤子携带费米零模，使最低能态空间成为 $\mathbf{L}^2(\mathcal{O}^+_m)\otimes S$（$n$ 为偶数时）。在 $3+1$ 维 $SO(3)$ 规范理论中孤子即磁单极子，其无穷远处第一 Chern 类 $c_1$ 必为偶数；加入 $\theta$ 角后正则动量被移动，于是 $Q_e=\frac{\theta}{2\pi}Q_m\bmod\mathbb{Z}$，晶格在 $\theta\to\theta+2\pi$ 下的单值变换给出 Witten 效应。
\end{keytakeaways}

% ===== 原文 p.129 =====
\subsection{什么是孤子？}
% 原文 §9.1：What is a soliton?

在经典数学物理中，所谓孤子（soliton）通常指非线性偏微分方程 $u_t=F(u,u_x,\dots)$ 的「行波」（traveling wave）解，即形如 $u(x,t)=f(x-vt)$ 的解。
孤子在可积系统理论中起着非常重要的作用：在该理论中（在某些情形下）任何解都可由以不同速度运动的孤子的叠加来逼近。
因此，可积系统理论有时也被称为孤子理论。

今天我们感兴趣的，则是场论中出现的孤子（作为经典场方程的行波解），并且主要是它们在场的量子化中所扮演的角色。
这与可积系统理论中的观点不同。
特别地，我们并不声称孤子之间存在非线性叠加；我们所考察的模型一般说来也不是精确可积的。

我们将考虑闵可夫斯基空间上庞加莱（Poincaré）不变的场论中的孤子。
对某个具体的场论，我们把孤子理解为场方程的一个行波解（即依赖于 $x-vt$ 的解），它在空间中是局域的，并且能量有限。
由庞加莱不变性，总可以假设 $v=0$，即该解与时间无关。
我们最关心的是这样的孤子：它在相应的同伦类中给出能量的全局极小。

\textbf{注（Remark）.}
把孤子与瞬子（instanton）区分开来十分重要。
瞬子局域于欧氏时空之中（即只存在一个瞬刻），其作用量有限；而孤子局域于（闵可夫斯基时空的）空间之中，永远存在，能量有限。
二者的物理应用完全不同。今天我们关心的是孤子。

% [Tutor Note: 「在相应同伦类中给出能量全局极小」这一限制并非多余。以 $1+1$ 维 $\phi^4$ 模型为例，分支 $X_{+-}$ 上能量泛函的临界点只有 kink（连同其平移与 boost），它自动是极小；但在分量更多的模型（$CP^{n}$ 模型、非阿贝尔规范理论）中，同一同伦类内可以有多个临界点，只有满足一阶 Bogomolny 方程的才是全局极小，其余是鞍点。作者在第 1249 页假设极小值非退化（横截方向二阶导数为正），正是为了保证孤子附近的态可按简谐振子量子化、而不会混入其它临界点的贡献。]

% [Tutor Note: 「由庞加莱不变性总可设 $v=0$」的准确含义是：若 $\phi(x,t)$ 是解，则任一 boost $\Lambda$ 给出新解 $\phi(\Lambda^{-1}(x,t))$，其能量为 $E\to\gamma E$。因此「在质心系中取能量极小」与「在任意惯性系中取能量极小」等价，孤子的静止系总是可以选定的；这也解释了后文 $P$-轨道 $\mathcal{O}=P/P_s$ 为何要以静止解的稳定化子 $P_s$ 作商。]

\subsection{孤子与经典解空间的连通分支}
% 原文 §9.2：Solitons and components of the space of classical solutions

经典地看，孤子的存在与「有限能量的经典解空间具有不同的连通分支」这件事密切相关。
在不包含零能量解的连通分支中，能量的极小值往往在某个孤子上达到。

% ===== 原文 p.130 =====
作为例子，可以回想第 1 讲中讨论过的情形：势为 $U=(\phi^2-a^2)^2$ 的二维标量场论。
此时，有限能量的、与时间无关的经典解空间有 4 个连通分支：$X_{++}$、$X_{-+}$、$X_{+-}$、$X_{--}$，其中
\[
X_{+-}=\{\phi:\ \phi(-\infty)=a,\ \ \phi(\infty)=-a\},
\]
其余符号取法类似。
在其中两个分支 $X_{+-}$ 与 $X_{-+}$ 上，能量严格为正，且其极小值分别在某个孤子解上达到。
这些解分别形如 $\phi=f(x)$ 与 $\phi=-f(x)$，其中 $f(x)$ 是势 $-U$ 所对应的牛顿方程 $f''=U'(f)$ 的、满足边界条件「在 $-\infty$ 处取值 $a$、在 $\infty$ 处取值 $-a$」的解。
证明这样的解存在、且在 $x$ 的平移下唯一，是初等的。

% [Tutor Note: 「$f''=U'(f)$，对应势 $-U$」这一句把两件事说成了一句话，容易读错：能量泛函是 $E=\int dx\,[\tfrac12(f')^2+U(f)]$，若把 $x$ 当作「时间」，则 $f(x)$ 满足粒子在势 $V=-U$ 中的牛顿方程 $f''=-V'(f)=U'(f)$。kink 相当于粒子从 $-U$ 的一个极大点「滚」到另一个极大点、且两端速度为零的轨道，能量守恒给出 $\tfrac12(f')^2=U(f)$，于是 $E=\int dx\,(f')^2$——这正是饱和了的 Bogomolny 界，也解释了为何 kink 是该分支上的全局极小而非鞍点。]

注意，孤子并不在庞加莱群下不变。
但人们常遇到在空间中绕某个「质心」具有旋转对称性的孤子。
在标量场论中，这意味着该孤子的对称群是 $P_s=SO(d-1)\times\mathbb{R}$。
由于庞加莱群是 $P=SO(d-1,1)\ltimes\mathbb{R}^{d-1,1}$，孤子在解空间中的 $P$-轨道 $\mathcal{O}$（即商 $P/P_s$）便是上半个双叶双曲面上的一个 $(d-1)$ 维向量丛。
以下几条容易验证。

% [Tutor Note: 作者称 $P/P_s$ 是「$(d-1)$ 维向量丛」，数一下维数便可确认：$\dim P=\frac{d(d-1)}{2}+d$，$\dim P_s=\frac{(d-1)(d-2)}{2}+1$，故 $\dim(P/P_s)=2(d-1)$，恰为底流形（双曲面，$(d-1)$ 维）的维数加上纤维维数 $(d-1)$；而 $T^*\mathcal{O}^+_m$ 的维数同样是 $2(d-1)$，这正是下面第 3 条中辛微分同胚得以存在的前提。]

\begin{enumerate}
\item 该丛自然同构于双曲面的余切丛。
\item 解空间上的辛结构限制到 $\mathcal{O}$ 上非退化，从而在 $\mathcal{O}$ 上定义了一个辛结构。
\item （辛形式的归一化）设 $m$ 为孤子的质量（即质心系中的能量）。
  则存在一个 $P$-等变的辛微分同胚
  \[
  \mathcal{O}\to T^*\mathcal{O}^+_m,
  \]
  其中 $\mathcal{O}^+_m$ 是 $\mathbb{R}^{d-1,1}$ 中双曲面 $x^2=-m^2$ 的上半部分。
\end{enumerate}

\subsection{孤子与量子化}
% 原文 §9.3：Solitons and quantization

在上述例子这类「经典解空间有若干不相连分支」的情形中（此例中即 $X_{\pm\pm}$），人们预期量子 Hilbert 空间是若干子空间的直和，每个分支对应一个子空间。
若 $X$ 是经典解空间的一个分支，且其中含有庞加莱不变的极小能量解，则可以预期 $X$ 的量子化将给出一个含有庞加莱不变态的 Hilbert 空间。
在上面的例子中，$X_{++}$ 与 $X_{--}$ 含有这样的庞加莱不变解（$\phi=a$ 或 $-a$）。
但 $X_{+-}$ 与 $X_{-+}$ 没有；在这些分支上，极小能量解是孤子或反孤子。
一般地，在不含庞加莱不变的最小能量解的分支中，最小能量解总是某种孤子。
因此我们把这样的分支称为孤子型（solitonic）分支。

设 $X_s$ 是一个孤子型分支。
那么，考虑孤子解的一切庞加莱变换，我们如上得到一个辛嵌入
\[
T^*\mathcal{O}^+_m\to X_s,
\]
其中 $m$ 是孤子解 $s$ 的质量。
我们假定能量在此映射的像上的极小值是非退化的，意思是：不存在能量为 $m$ 的其它解，且能量沿横截于该像的方向的二阶导数为正。

在此情形下，在弱耦合区域我们应当预期：

% ===== 原文 p.131 =====
\begin{enumerate}
\item 量子理论 Hilbert 空间 $\mathcal{H}$ 的分支 $\mathcal{H}_s$ 没有真空（因为经典上没有庞加莱不变解）。
\item $\mathcal{H}_s$ 上的哈密顿量 $H$ 满足不等式 $H\ge m'$，其中 $m'$ 是某个正的质量参数，且当耦合趋于零时 $m'\to m$（一般说来，应当预期 $m$ 会得到量子修正）。
  若理论有质量间隙，则应预期：当耦合趋于零时质量收敛到 $m$ 的那些态，构成 $T^*\mathcal{O}^+_m$ 的一个量子化，即 $P$ 的形如 $\mathbf{L}^2(\mathcal{O}^+_m)$ 的不可约表示。
  由于二阶导数为正，还应预期 $m$ 附近的谱是离散的。
\end{enumerate}

注意，若理论的（作用量或）能量泛函乘以一个常数 $C$，则孤子的质量也乘以 $C$，而通常粒子的质量不变。
回想一下，在这种情形下耦合常数按 $C^{-1}$ 标度。
这意味着，在半经典（弱耦合）近似（$C\to\infty$）下，孤子比通常粒子重得多。
因此，孤子不可能在微扰论中被看到：含孤子的中间态对关联函数的贡献在弱耦合极限下是指数级小的。
这正是我们此前能够在完全不讨论孤子的情况下走到这一步的原因！

% [Tutor Note: 「$C\to\infty$ 即弱耦合」的依据是把 $C$ 吸收进 $\hbar$：路径积分为 $\int\mathcal{D}\phi\,e^{\frac{i}{\hbar}C S_0}$，故展开参数是 $\hbar/C$，令 $\hbar=1$ 后即为 $1/C$，与作者所说的「耦合按 $C^{-1}$ 标度」一致。孤子质量 $M=C\,M_{\text{cl}}$ 随 $C$ 线性增长，微扰粒子质量则与 $C$ 无关；于是孤子中间态给出 $\sim e^{-M\Delta t}=e^{-C M_{\text{cl}}\Delta t}$ 的压低。注意这里说的是「指数小」而非「幂次小」，故任何有限阶微扰展开都看不见孤子。]

% [Tutor Note: 第 2 条中 $H\ge m'$ 与「无真空」是同一拓扑事实的两种表述：孤子分支的经典解在空间无穷远处取非平凡边界条件，任何有限能量的场组态都无法连续变形到 $\phi\equiv a$，故量子化后该分支的 Hilbert 空间里不存在 $P$-不变的基态。$m'$ 与 $m$ 的差别正是 $1/C$ 阶的量子修正（零点能、Casimir 型贡献），作者用 $m'\to m$ 把它一笔带过。]

\subsection{含费米子的理论中的孤子}
% 原文 §9.4：Solitons in theories with fermions

在含费米子的理论中，孤子在 $P$ 下的轨道往往不是解空间相应连通分支中最低能态的全体。
例如，若模型是超对称的，则先验地显然该轨道不是最低能量解的全体：最低能态的空间是超庞加莱群的轨道，而不只是庞加莱群的轨道。
即便没有超对称，费米子也常常在孤子的量子化中起重要作用。

让我们看一个这样的例子。考虑二维的、含一个标量与一对费米子的理论：
\begin{equation}
  \tag{9.1}
  \label{eq:9-1}
  \mathcal{L}=\frac{1}{2\lambda}\int d^2x\bigl(|\partial\phi|^2+(\phi^2-a^2)^2\bigr)
  +i\int d^2x\bigl(\psi_+\partial_-\psi_++\psi_-\partial_+\psi_--g\phi\psi_+\psi_-\bigr).
\end{equation}

\textbf{注 1（Remark 1）.}
对 $g$ 的某个适当取值，此模型是超对称的；但我们并不限于这个取值。

\textbf{注 2（Remark 2）.}
此模型具有手征对称性 $\phi\to-\phi$、$\psi_\pm\to\pm\psi_\pm$，它禁止费米子的质量项 $m\psi_+\psi_-$。

考虑理论玻色部分的孤子 $\phi(x,t)=f(x)$（在 $\phi=a$ 与 $\phi=-a$ 之间插值）。
此解的费米子推广是有限能量的函数 $\psi(x)=\begin{pmatrix}\psi_+(x)\\ \psi_-(x)\end{pmatrix}$，满足 Euler–Lagrange 方程
\begin{equation}
  \tag{9.2}
  \label{eq:9-2}
  \left[\partial_x\psi+\begin{pmatrix}0&-gf(x)\\-gf(x)&0\end{pmatrix}\right]\psi=0.
\end{equation}
于是
\begin{equation}
  \tag{9.3}
  \label{eq:9-3}
  \psi=\begin{pmatrix}\varepsilon\\ \varepsilon\end{pmatrix}
  \exp\Bigl(g\int_0^x f(y)\,dy\Bigr)
  +\begin{pmatrix}\varepsilon'\\ -\varepsilon'\end{pmatrix}
  \exp\Bigl(-g\int_0^x f(y)\,dy\Bigr),
\end{equation}

% [Tutor Note: (9.2) 的记号略有滥用：方括号内是作用在二分量 $\psi$ 上的算符 $\partial_x+\begin{pmatrix}0&-gf\\-gf&0\end{pmatrix}$，而 $\psi$ 同时出现在方括号内外；按字面读作 $[\partial_x+M(x)]\psi=0$ 即可。(9.3) 就是该常微分方程的通解：矩阵 $M$ 的两个本征向量 $(1,1)$ 与 $(1,-1)$ 分别对应本征值 $\mp gf$，故指数分别是 $\pm g\int_0^x f$。]

% ===== 原文 p.132 =====
其中 $\varepsilon,\varepsilon'$ 是奇变量。
容易看出，对孤子（$f>0$ 当 $x\to\infty$，$f<0$ 当 $x\to-\infty$）而言，恰好 $\varepsilon=0$ 的那些解属于 $\mathbf{L}^2$，于是有限能量解的空间是一维的。
对反孤子同样，恰好 $\varepsilon'=0$ 的那些解属于 $\mathbf{L}^2$。
在真空分支中（$f$ 在 $x\to\pm\infty$ 时同号），没有 $\mathbf{L}^2$ 解。

% [Tutor Note: 判断哪个指数被 $\mathbf{L}^2$ 排除，只看 $x\to\pm\infty$ 的渐近行为即可：$f(\pm\infty)=\pm a$，故 $\exp(\pm g\int_0^x f)$ 分别按 $e^{\pm gax}$ 增长或衰减。第一项在 $x\to+\infty$ 发散，必须取 $\varepsilon=0$；第二项在两端都衰减，于是只剩一个（由奇参数 $\varepsilon'$ 标记的）归一化解——这就是「恰有一个费米零模」。零模的存在依赖于 $f$ 的拓扑（kink）；真空分支中两端同号，两个指数在某一端都发散，故无 $\mathbf{L}^2$ 解。]

于是，考察孤子分支，我们便得到费米子的唯一一个零能 $\mathbf{L}^2$ 解。
我们可以把模型稍作推广，使之不含单独一对费米场 $\psi_+,\psi_-$，而含 $n$ 对 $\psi_{\pm i}$，$i=1,\dots,n$，每一对都通过与方程 (9.1) 中相同的耦合与 $\phi$ 相互作用。
于是拉氏量的费米子部分为
\[
\sum_{i=1}^{n}i\int d^2x\bigl(\psi_{+i}\partial_-\psi_{+i}
+\psi_{-i}\partial_+\psi_{-i}-g\phi\psi_{+i}\psi_{-i}\bigr).
\]
费米子上按自然方式作用着 $SO(n)$ 对称性。
此时有 $n$ 个费米子解，故最低能量解的庞加莱轨道不是 $T^*\mathcal{O}^+_m$，而是超流形 $T^*\mathcal{O}^+_m\times\mathbb{R}^{0|n}$。

因此，若 $n$ 为偶数，量子理论中 Hilbert 空间相应分支里的最低能态空间是 $\mathbf{L}^2(\mathcal{O}^+_m)\otimes S$，其中 $S=S_++S_-$ 是 $Spin(n)$ 的自旋表示。
$n$ 为奇数时描述略为微妙，今天我们不去分析这种情形。（关于其要点的说明，见第 12 讲第 3 部分的末尾。）

% [Tutor Note: $n$ 为偶数时 $Spin(n)$ 有两个不等价的半自旋表示 $S_\pm$，费米数的 $\mathbb{Z}_2$ 分次与 $S=S_++S_-$ 的分解相容，故基态空间 $\mathbf{L}^2(\mathcal{O}^+_m)\otimes S$ 有良好定义。$n$ 为奇数时 $Spin(n)$ 只有一个自旋表示，$Cl(n)$ 的不可约模不能自然地按费米数分半，孤子基态的费米奇偶性因而没有先验定义——这正是作者称「描述更微妙」并指向第 12 讲第 3 部分的原因。]

\subsection{$2+1$ 维与 $3+1$ 维中的孤子}
% 原文 §9.5：Solitons in 2+1 and 3+1 dimensions

尽管瞬子与孤子有区别，二者之间仍有联系。
即：$d-1$ 维欧氏场论中的一个瞬子，往往给出同一理论 $d$ 维闵可夫斯基版本中的一个孤子。
例如，考虑带复标量的 $2+1$ 维 $U(1)$ 规范理论：
\begin{equation}
  \tag{9.4}
  \label{eq:9-4}
  \mathcal{L}=\int d^3x\left(\frac{F_A^2}{4e^2}
  +\frac{1}{\lambda}\bigl(|d_A\phi|^2+(|\phi|^2-a^2)^2\bigr)\right).
\end{equation}
在第 7 讲我们看到，此理论的二维版本有 Chern 类为 $1$ 和 $-1$ 的瞬子。
在 $2+1$ 维理论中，这些瞬子变成孤子，并如上所述支配 Hilbert 空间相应分支中的最低能模式。

现在考虑 $3+1$ 维的非阿贝尔规范理论。
即考虑带三维表示中玻色子的 $SO(3)$ 规范理论，其拉氏量为
\begin{equation}
  \tag{9.5}
  \label{eq:9-5}
  \mathcal{L}=\int d^4x\left(\frac{\operatorname{Tr}F_A^2}{4e^2}
  +\frac{1}{2\lambda}\bigl(|d_A\phi|^2+(|\phi|^2-a^2)^2\bigr)\right).
\end{equation}
% [Tutor Note: 注意 (9.4) 与 (9.5) 中标量势项的系数不同：$1/\lambda$ 对 $1/(2\lambda)$（规范项则都是 $F_A^2/4e^2$）。这不是笔误：$d=3$ 时 $[\phi]=1/2$、$[\lambda]=-1$，$d=4$ 时 $[\phi]=1$、$[\lambda]=0$，两处的 $\lambda$ 是量纲不同的参数，相对因子 $2$ 可完全吸收进 $\lambda$ 的定义。但若要把两式放在一起比较（例如核对单极子质量公式 $M\sim4\pi a/e$ 的系数），就必须先统一约定，否则会凭空差一个因子 $2$。]
我们可以把 $A$ 看作秩三实向量丛 $E$ 上的联络，把 $\phi$ 看作 $E$ 的截面。
我们在第 8 讲中考虑过此理论的三维版本。
为方便起见，我们把时空 $\mathbb{R}^3$ 与规范群的李代数、以及 $\phi$ 的取值空间等同起来（李代数中的括号就是叉积）。
这使我们能把标量场 $\phi$ 与无穷小规范变换写成 $\mathbb{R}^3$ 上的向量场。

% ===== 原文 p.133 =====
我们发现，在三维情形此理论有一个瞬子，其中 $\phi$ 形如
\[
\vec\phi=\frac{\vec x}{r}f(r),\qquad r=|\vec x|
\]
（如我们已经说明的，我们把时空与 $\phi$ 的取值空间等同起来，故 $\vec\phi$ 与位置 $\vec x$ 都是三维向量）。
在 4 维中，这个瞬子将变成一个孤子。
这类孤子称为磁单极子（magnetic monopole），因为它们相对于低能未破缺的 $U(1)$ 规范理论确实带有经典意义上的磁荷。

为更好地理解这一点，回想：任何有限能量的场组态 $(A,\phi)$ 都定义一个整值拓扑不变量——「无穷远处的第一 Chern 类 $c_1$」。
的确，为使能量有限，必须在无穷远处有 $|\phi|=a$，于是 $\phi$ 定义了从 $\mathbb{R}^3$ 中无穷远处的球面到半径为 $a$ 的球面的一个映射，而 $c_1$ 是这个映射的度的两倍。

$c_1$ 的另一个定义（它或许更能说明该不变量为何叫 $c_1$、以及在何种意义上 $c_1$ 是上述映射的度的两倍）如下：截面 $\phi$ 在无穷远处把我们的三维向量丛 $E$ 分裂为直和 $\mathbb{R}\phi\oplus\phi^{\perp}$，即一个一维与一个二维向量丛。
（这里 $\mathbb{R}\phi$ 是 $E$ 中由与 $\phi$ 成比例的向量构成的子丛，$\phi^{\perp}$ 是正交补。）
于是 $\phi$ 在无穷远处把结构群约化为 $SO(2)=U(1)$。
数 $c_1$ 就是 $U(1)$ 丛 $\phi^{\perp}$ 限制到 $\mathbb{R}^3$ 中无穷远 2-球面上所得的第一 Chern 类。
例如，上面讨论的孤子组态有 $c_1=2$，因为此时 $SO(2)$ 丛正是无穷远处球面的切丛，而球面的 Euler 示性数为 $2$。
在本模型中，在无穷远处测得的 $c_1$ 总是偶数，故 $2$ 是最小的可能正值。
其为偶数的理由是：由于 $\mathbb{R}^3$ 可缩，$E$ 的第二 Stiefel–Whitney 类为零，于是我们可以过渡到 $SO(3)$ 的双覆盖 $SU(2)$，并考虑与 $E$ 相伴的秩二复向量丛 $W$。
Higgs 场 $\phi$ 在无穷远处把 $W$ 分裂为 $\mathcal{T}\oplus\mathcal{T}^{-1}$，其中 $\mathcal{T}$ 是某个线丛；$\mathcal{T}$ 的第一 Chern 类等于用 $SO(3)$ 丛定义的 $c_1$ 的一半，故后者总是偶数。
当然，无穷远处的第一 Chern 类在第 8 讲中也以大致相同的方式引入过。

% [Tutor Note: 「$c_1$ 是映射 $\phi/|\phi|:S^2_\infty\to S^2_a$ 的度的两倍」与「上述孤子组态有 $c_1=2$」必须一致地理解，否则容易以为差了因子 2：标准 't Hooft–Polyakov 单极子取 $\phi^a=\hat x^af(r)$，$\phi/|\phi|$ 是恒等映射（度 $1$），而 $\phi^{\perp}$ 在 $S^2_\infty$ 上的限制恰为 $TS^2$，$\int_{S^2}c_1(TS^2)=\chi(S^2)=2$，故 $c_1=2$、磁荷 $c_1/2=1$。因子 $2$ 正是「$SO(3)$ 丛的 $c_1$ 等于 $SU(2)$ 提升下线丛 $c_1$ 的两倍」这一事实的另一面，也是本模型中 $c_1$ 必为偶数的根源。]

% [Tutor Note: 「$\mathbb{R}^3$ 可缩故 $w_2(E)=0$，于是可过渡到双覆盖 $SU(2)$」这一步省略了：$w_2\in H^2(\mathbb{R}^3;\mathbb{Z}_2)=0$ 是 $SO(3)$ 丛提升为 $SU(2)$ 丛的障碍类，障碍为零故提升存在；而 $SU(2)$ 的秩二复表示 $W$ 上，$\phi$ 在无穷远处给出 $W=\mathcal{T}\oplus\mathcal{T}^{-1}$ 的分解，$\mathcal{T}$ 与 $\mathcal{T}^{-1}$ 的 $c_1$ 互为相反数，故 $c_1(W)=0$，这与 $W$ 来自 $SU(2)$ 丛相容。]

此理论的孤子解被称为磁单极子，理由如下。
在无穷远处（或在真空中），规范对称性通过我们本学期第 2 讲所研究的 Higgs 机制自发破缺，从 $SO(3)$ 破缺到 $SO(2)=U(1)$。
于是低能理论看起来像一个 $U(1)$ 阿贝尔规范理论。
无穷远处的第一 Chern 类度量的是「磁场」（低能 $U(1)$ 理论的曲率 $F/2\pi$）在无穷远处球面上的积分；在经典物理中，这个积分被称为总磁荷。（我们将在下文 (9.9) 式给出磁荷的形式定义。）
人们通常把 $c_1/2$ 称为该解的磁荷。

现在讨论这样一个事实对孤子的一项推论：无穷远处规范群的未破缺子群是 $U(1)$。
这个 $U(1)$ 是真空、或无穷远处场的对称性；但它并不保持完整的孤子（磁单极子）解不变。
作用在某个特定解上，它会产生一个新的解。
为理解这个作用，让我们考察上面讨论过的、$c_1=2$ 的孤子 $(A,\phi)$。
我们把 $U(1)$ 群的无穷小生成元表示为球对称的规范对称性：$\varepsilon=\frac{\vec x}{r}g(r)$，其中 $g$ 是某个函数。
若函数 $g(r)$ 满足 $r\to0$ 时 $g(r)\sim c_0r$、以及 $g(+\infty)=c$，则此式定义了规范变换李代数中一个光滑元素，它在无穷远处是「常数」的（相对于由标量场 $\phi$ 所定义的结构群约化）。

% ===== 原文 p.134 =====
若愿意，我们可以取 $\varepsilon=\phi/a$，这是一个完全自然、且满足全部条件的选择。
于是 $\varepsilon$ 的作用为
\begin{equation}
  \tag{9.6}
  \label{eq:9-6}
  A\to A-d_A\varepsilon,\qquad \phi\to\phi.
\end{equation}
显然 $d_A\varepsilon$ 不可能恒为零，因为孤子的联络 $A$ 并不具有整体约化到 $U(1)$ 的结构群；因此不存在协变常数的 $\varepsilon$。
所以，位于无穷远处未破缺群中的那些规范变换，其作用是非平凡的。
现在考虑作规范变换 $\exp(c\varepsilon)$。
若 $c\in2\pi\mathbb{Z}$，则该规范变换在无穷远处趋于恒等。
在量子化时，我们把相差一个「在无穷远处趋于恒等」的规范变换的两个场视为同一，故对这样的 $c$，应认为 $\exp(c\varepsilon)$ 平凡地作用。
模去在无穷远处等于 $1$ 的规范变换，那些保持无穷远处场不变的规范变换生成未破缺的 $U(1)$，而这个群自由地作用在孤子解的空间上。

% [Tutor Note: 取 $\varepsilon=\phi/a$ 时，由 $\phi=\frac{\vec x}{r}f(r)$ 得 $g(r)=f(r)/a$；光滑性要求 $r\to0$ 时 $f(r)\sim c_0'r$，故 $g(r)\sim(c_0'/a)r$，又 $f(\infty)=a$ 给出 $g(+\infty)=1$，两条要求恰好满足。这也说明「常数 $c$」的量子化条件为何是 $c\in2\pi\mathbb{Z}$：$\varepsilon=\phi/a$ 在无穷远处取值于单位球面，而 $so(3)$ 中绕单位轴转 $2\pi$ 回到恒等，故 $\exp(c\varepsilon)|_{\infty}=1$ 当且仅当 $c$ 是 $2\pi$ 的整数倍。]

% [Tutor Note: 「依赖时间的孤子空间是 $T^*\mathcal{O}\times T^*S^1$」这一步是作者跳过的：静态孤子的模空间是 $\mathcal{O}\times S^1$（质心位置与 $\alpha$ 角），而场论中任一静态解的速度自由度总由「模空间坐标 + 其共轭动量」参数化，故依赖时间解的模空间是余切丛。平移给出 $T^*\mathbb{R}^3\subset T^*\mathcal{O}$，而 $\alpha=\alpha_0+st$ 给出 $T^*S^1$ 的纤维方向；注意 $s$ 只是 $\alpha$ 的「速度」，与 $\alpha$ 共轭的动量才是后文的电场荷 $Q_e$，二者差一个由辛形式归一化决定的常数。]

同时允许空间平移，上述论证表明：磁荷为 $1$ 的、依赖时间的孤子构成的空间至少是 $\mathbb{R}^3\times S^1$（$\mathbb{R}^3$ 与 $S^1$ 上都没有典则的零点）。
可以证明，事实上它恰好就是如此。
我们把 $S^1$ 因子上的圆坐标记作 $\alpha$。

利用庞加莱群变换，我们可以生成依赖时间、以常速传播的解。
我们也可以作依赖时间的规范变换，这会使 $\alpha$ 坐标依赖时间，即 $\alpha=\alpha_0+st$，$s\in\mathbb{R}$。
结果，依赖时间的孤子构成的空间是乘积 $T^*\mathcal{O}\times T^*S^1$。

\subsection{带 $\theta$ 角的 $3+1$ 维理论}
% 原文 §9.6：The 3+1-dimensional theory with the θ-angle

考虑上一节 $3+1$ 维理论并加入 $\theta$ 角，即在我们（闵可夫斯基）的拉氏量中添上一项
\begin{equation}
  \tag{9.7}
  \label{eq:9-7}
  -\frac{\theta}{16\pi^2}\int\operatorname{Tr}(F\wedge F).
\end{equation}
我们回想，第二 Chern 类 $c_2$ 由积分 $\frac{1}{8\pi^2}\int\operatorname{Tr}(F\wedge F)$ 给出。
因此，经典地看，项 (9.7) 只是数出丛的第二 Chern 类。

现在考虑我们理论中的三个算符：
\begin{enumerate}
\item 未破缺 $U(1)$ 的生成元，我们把它写作 $\frac{\partial}{\partial\alpha}$，其中 $\alpha$ 是单极子解空间圆因子上的角参数。
\item 电荷
  \begin{equation}
    \tag{9.8}
    \label{eq:9-8}
    Q_e=\frac{1}{ae^2}\int_{\mathbb{R}^3}\operatorname{Tr}(d_A\phi\wedge *F_A).
  \end{equation}
\item 磁荷
  \begin{equation}
    \tag{9.9}
    \label{eq:9-9}
    Q_m=\frac{1}{4\pi a}\int_{\mathbb{R}^3}\operatorname{Tr}(d_A\phi\wedge F_A).
  \end{equation}
\end{enumerate}

% ===== 原文 p.135 =====
通过分部积分，并利用 $F$ 的运动方程与 Bianchi 恒等式，$Q_e$ 与 $Q_m$ 都可以写成无穷远处的面积分；这些面积分可以写成 $\ast F/2\pi$ 或 $F/2\pi$ 在无穷远处球面上的积分，其中 $F$ 是未破缺 $U(1)$ 的曲率。
于是，这些定义在低能理论中归结为第 8 讲引入的 $Q_e$ 与 $Q_m$。
特别地，$Q_m$ 就是磁荷、即无穷远处的第一 Chern 类，从而是经典解空间上的局部常函数。
就我们当前的目的而言，与其分部积分，不如把这些公式保留为这里给出的形式。

对由场量子化而来、并进入微扰论的普通粒子，$Q_m=0$（因为这些粒子是通过量子化 $c_1=0$ 的那部分经典解空间得到的），且 $Q_e=-i\frac{\partial}{\partial\alpha}$。（若此陈述在前面的讲次中不熟悉，稍后就会给出推导。）
于是在微扰论中，$Q_e$ 具有整本征值。

然而，让我们看看非微扰地会发生什么，即当我们把 $c_1\neq0$ 的场组态也考虑进来时。

我们在 $A_0=0$ 的规范下计算；把 $dA/dt$ 记作 $\dot A$。
经典地有
\[
\frac{\partial}{\partial\alpha}
=a^{-1}\int\operatorname{Tr}\Bigl(d_A\phi\,\frac{\delta}{\delta\dot A}\Bigr).
\]
理由是：如上面提到的，我们可以认为未破缺的 $U(1)$ 由李代数取值函数 $\varepsilon=-\phi/a$ 生成。
由该函数生成的规范变换保持 $\phi$ 不变，并在一阶上把 $A$ 变换为 $\delta A=d_A\phi/a$。
这正是由 $a^{-1}\int\operatorname{Tr}\bigl(d_A\phi\,\frac{\delta}{\delta\dot A}\bigr)$ 生成的变换，说明该算符就是 $\frac{\partial}{\partial\alpha}$。

因此，算符 $-i\frac{\partial}{\partial\alpha}$ 是流 $J=\operatorname{Tr}(d_A\phi\wedge\pi_A)$ 的荷，其中 $\pi_A$ 是 $A$ 的共轭（动量）变量。

现在计算 $\pi_A$：
\begin{equation}
  \tag{9.10}
  \label{eq:9-10}
  \pi_A=\frac{\delta\mathcal{L}}{\delta\dot A}
  =\frac{\ast F_A}{e^2}-\frac{\theta}{8\pi^2}F_A.
\end{equation}
把它与 $\partial/\partial\alpha$ 的表达式以及 $Q_e$、$Q_m$ 的定义结合起来，得到
\begin{equation}
  \tag{9.11}
  \label{eq:9-11}
  -i\frac{\partial}{\partial\alpha}=Q_e-\frac{\theta}{2\pi}Q_m.
\end{equation}
在微扰论中 $Q_m=0$，故如上所述 $Q_e=-i\partial/\partial\alpha$。

由于 $-i\frac{\partial}{\partial\alpha}$ 的谱是整数，我们得到 $Q_e=\frac{\theta}{2\pi}Q_m\bmod\mathbb{Z}$。
但 $Q_m$ 也是整数，因为它在场的每个分支上都是整数。
于是，对孤子态而言 $Q_e$ 不必是整数。
换言之，存在这样的态（在半经典极限下能量 $\sim1/\hbar$，非常高），其上的 $Q_e$ 不是整数。
所以，在没有磁单极子时电荷是整数的，但在有磁单极子时则未必。
大多数物理学家猜测，自然界中存在带非整数、且很可能是无理数电荷的磁单极子，其机制本质上就是上面这个。

% [Tutor Note: (9.10) 中 $\theta$ 项的系数是 $\theta/8\pi^2$ 而非 $\theta/16\pi^2$，这一步容易算错：对拓扑项 $\int\operatorname{Tr}(F\wedge F)$ 关于 $\dot A$ 变分时，$\varepsilon^{\mu\nu\rho\sigma}$ 把 $F\wedge F$ 中的一个 $F$ 收缩成 $\ast F$ 并给出因子 $2$，正好抵消 (9.7) 中的 $1/16\pi^2$。其物理含义是：$\theta$ 项在经典层面是全导数、不改变场方程，却移动了正则动量 $\pi_A$，从而移动了电荷谱——这正是「Witten 效应」的核心。]

% [Tutor Note: 原文称单值变换矩阵在基 $Q_m,Q_e$ 下为 $\begin{pmatrix}1&1\\0&1\end{pmatrix}$；但若把 $(Q_m,Q_e)^{T}$ 当作列向量，实现 $Q_m\to Q_m$、$Q_e\to Q_e+Q_m$ 的矩阵应是 $\begin{pmatrix}1&0\\1&1\end{pmatrix}$，所写矩阵恰为其转置。因此这里疑为基序笔误（所写矩阵相当于在基序 $(Q_e,Q_m)$ 下作用，或让矩阵作用在行向量/对偶格上）。物理结论——$\theta\to\theta+2\pi$ 时 $Q_e$ 平移 $Q_m$——不受影响。]

% [Tutor Note: 「能量 $\sim1/\hbar$」指的是这些态本身是孤子（磁单极子）激发：孤子质量 $\sim1/\hbar$，在 $\hbar\to0$ 时趋于无穷，故它们不会出现在微扰谱中；而 $Q_e\to Q_e+Q_m$ 只是在此态上叠加 $\theta$ 依赖的电场，不改变其能量量级。这也解释了为何微扰分支（$Q_m=0$）中电荷仍取整数。]

算符 $Q_e,Q_m$ 对易，因为经典上 $Q_m$ 是局部常函数。
$Q_m$ 与 $-i\partial/\partial\alpha$ 都取整数值，结合 (9.11) 便表明 $Q_e,Q_m$ 的联合谱是 $\mathbb{R}^2$ 中的一个晶格。
若考虑由 $\theta\in S^1$ 参数化的理论族，则该晶格绕此圆一周的单值变换由基 $Q_m,Q_e$ 下的矩阵 $\begin{pmatrix}1&1\\0&1\end{pmatrix}$ 给出。
这一断言是 (9.11) 的推论：它表明，当绝热地把 $\theta$ 增加 $2\pi$ 时，$Q_e$ 增加 $Q_m$。

% ===== 原文 p.136 =====
我们刚刚完成的计算，在后文讨论 Donaldson 理论时会很重要，它与第 8 讲末尾在阿贝尔理论中所作的计算密切相关。
利用显式公式 (8.78)，我们在第 8 讲注意到：在 $\theta\to\theta+2\pi$ 之下，通过任一二维曲面 $\Sigma$ 的电通量与磁通量 $Q_e(\Sigma)$、$Q_m(\Sigma)$ 按 $Q_e\to Q_e+Q_m$ 变换。
在阿贝尔理论中，为使 $Q_e$ 与 $Q_m$ 不为零，我们必须在一个第二 Betti 数非零的三维流形上工作。
而在当前非阿贝尔情形的讨论中，我们可以简单地在 $\mathbb{R}^3$ 上工作，取 $\Sigma$ 为无穷远处的球面。
我们这里所称的 $Q_e$ 与 $Q_m$，就是 $Q_e(\Sigma)$ 与 $Q_m(\Sigma)$，在低能有效理论中、取 $\Sigma$ 为无穷远处球面时算得。

在第 8 讲中我们还看到（在自由理论中）如何把 $\theta\to\theta+2\pi$ 下的这个单值表示推广为 $SL(2,\mathbb{Z})$ 的作用，即如何构造由复参数 $\tau=\frac{\theta}{2\pi}+ir$ 参数化、模去模变换后的一族理论，使得单值表示恰是 $SL(2,\mathbb{Z})$ 在二维晶格上的标准作用。
要在四维非线性理论中做到类似的事情，是更困难的问题，我们只在学期末才会回到它。

% [Tutor Note: $\tau=\frac{\theta}{2\pi}+ir$ 的虚部 $r$ 与 $\theta$ 无关，它扮演的是「规范耦合常数」的角色（$r\sim1/e^2$）；$SL(2,\mathbb{Z})$ 的单值性来自两个生成元 $\tau\to\tau+1$（即 $\theta\to\theta+2\pi$，本讲已显式验证）与 $\tau\to-1/\tau$（强弱对偶，在四维渐近自由理论中只是一类对偶猜测，并非已证定理）。作者说「回到这一点要等到学期末」，指的是后面讲 Seiberg–Witten 理论时才会真正构造这种族。注意 $\theta$ 与 $r$ 不能独立延拓：模变换把 $\theta$ 与 $1/e^2$ 混合，这正是 Witten 效应在 $\tau$ 语言下的表述。]

\subsection{参考文献与交叉引用}

\begin{itemize}
\item 关于孤子的入门介绍，见 S. Coleman, \emph{Aspects of Symmetry}, Cambridge University Press, 1985 中 “Lumps” 一章。
\item 关于磁单极子电荷之 $\theta$ 依赖性的推导，见 E. Witten, \emph{Dyons of Charge $e\theta/2\pi$}, Phys. Lett. B \textbf{86} (1979), 283--287。
\item 超对称理论中的孤子（BPS 态）将在第 12 讲以及本系列后续各讲中讨论。3 维与 4 维中的孤子曾在第 6、7 讲中讨论过（作为 2 维与 3 维中的瞬子）。
\end{itemize}
